% ==============================================================================
% DECLARAÇÃO SOBRE O USO DE IA GENERATIVA E TECNOLOGIAS ASSISTIDAS POR IA
% ==============================================================================

\section*{Declaração sobre o uso de IA generativa}
\label{sec:declaracao-uso-ia}

Durante a elaboração deste Trabalho de Conclusão de Curso e a construção da plataforma DeployDock, foram utilizadas tecnologias assistidas por inteligência artificial generativa baseadas em modelos de linguagem de grande porte (LLMs). As ferramentas atuaram como copilotos no ciclo de vida de desenvolvimento de software, conforme pormenorizado na seção de Metodologia, participando de:
\begin{enumerate}
    \item Apoio na modelagem e síntese de expressões regulares heurísticas para o Web Application Firewall (WAF) contra ataques de Injeção de SQL, Cross-Site Scripting, Path Traversal e Execução Remota de Código;
    \item Elaboração de algoritmos de detecção multilinguagem para geração automatizada de receitas de contêineres (\textit{Auto-Dockerfile Zero-Config});
    \item Estruturação da suíte de 58 testes automatizados ponta a ponta (E2E) e rotinas assíncronas para a bateria de integridade transacional ACID do PostgreSQL 16;
    \item Suporte à redação técnica, formatação de tabelas e equações segundo as normas e padrões da ABNT e do template institucional.
\end{enumerate}

Todos os algoritmos, regras de segurança, código-fonte da aplicação e dados de testes empíricos foram minuciosamente inspecionados, revisados, validados e submetidos a testes práticos em ambiente real pelo autor, que assume integral e irrestrita responsabilidade científica e acadêmica por todo o conteúdo deste trabalho.
